%%%%%%%%%%%%%%%%%%%%%%%%%%%%%%%%%%%%%%%%%%%%%%%%%%%%%%%%%%
%% Template for small lab reports
%% Olivier Morel
%% inspired
%% template from reds institute for syslog
%%%%%%%%%%%%%%%%%%%%%%%%%%%%%%%%%%%%%%%%%%%%%%%%%%%%%%%%%%

% Paramétrage du papier utilisé
\documentclass[a4paper,12pt,oneside]{article}

% Ne pas toucher, nécessaire à la mise en forme
\usepackage[utf8x]{inputenc} 
\usepackage[T1]{fontenc} % Nouvelle norme pour codage des caractères
\usepackage{lmodern} % Nouvelle forme de la fonte ComputerModern
\usepackage[french]{babel} % Règles typographiques françaises, césures
\usepackage{graphicx} % Insertion images
\usepackage{float}  % allow to force the placement of picture with 'H'
\usepackage{epstopdf} % Convertir les images eps vers pdf
\usepackage[top=3cm, bottom=3cm, left=3cm, right=3cm, showframe=false]{geometry}
\usepackage{amsmath}
%\usepackage{hyperref}  % used to produce hyperlink when referencing but those links are marked by big red square on pdf softwares (in theory not when printed)
\usepackage{pdfpages}   % allow to include multipages pdf in document
\usepackage{color}  % allow to use colors
\usepackage{multirow}
\usepackage{multicol}
\usepackage{lastpage}
\usepackage{siunitx}

\epstopdfsetup{outdir=./Images/converted_to_pdf/} % Dossier de sortie des images converties
\pagenumbering{arabic}
\renewcommand{\familydefault}{\sfdefault}	% change font to helvetica
\sisetup{
	locale = FR,
	round-mode = figures,
	round-precision = 4,				% 4 digits in total
	scientific-notation = engineering,	% exponnant by multiple of 3
	per-mode            = symbol,
	per-symbol          = /,
	exponent-to-prefix  = true,
	output-decimal-marker = {.},		% decimal marker will always be printed as a '.'
	output-complex-root = j,			% imaginary unit will always be 'j'
	exponent-product = \cdot			% exponant product will be a dot
}

% fichier avec les constantes (à adapter pour les étudiants)
%%%%%%%  Constantes (à changer) %%%%%%%%

% add any number of students
\newcommand{\nomAuteur}{Etudiant1 \& Etudiant2}

\newcommand{\titre}{Titre du Laboratoire}

\newcommand{\nbrlabo}{X} % numéro de laboratoire

\newcommand{\prof}{P. Rof}

\newcommand{\assistant}{A. Sisstant}

\newcommand{\classe}{Ma classe}

\newcommand{\departement}{Technologies industrielles (TIN)}

\newcommand{\cours}{Mon Cours 1 (Mcou1)}


% Entête et pieds de page (A ADAPTER)
\pagestyle{empty}
\usepackage{fancyhdr}
\setlength{\headheight}{27.06pt}
\pagestyle{fancy}
% Entête :
\renewcommand{\headrulewidth}{1pt}
\fancyhead[L]{Labo\nbrlabo: \titre} 
\fancyhead[R]{Rapport de laboratoire} 
% Pied de page
\fancyfoot{} % clear all footer fields
\renewcommand{\footrulewidth}{1pt}
\fancyfoot[R]{Page \thepage / \pageref{LastPage}} 
\fancyfoot[L]{HEIG-VD | \nomAuteur }

\begin{document}	
	% première page
	%%%%%%%%%%%%%%%%%%%%%%%%%%%%%%%%%%%%%%%%%%%%%%%%%%%%%%%%
% Ce document ne doit pas être modifié
% Il faut modifier les informations depuis le fichier
% constantes.tex
%%%%%%%%%%%%%%%%%%%%%%%%%%%%%%%%%%%%%%%%%%%%%%%%%%%%%%%%



\begin{center}
\includegraphics[height=2cm]{./images/logo_heig.pdf}
\hfill \includegraphics[height=2cm]{./images/logo_hesso.pdf}
\end{center}

\vfill

\begin{center}
	\Large Laboratoire \nbrlabo:\\
	\huge \titre \\
\end{center}


\begin{center}
	\begin{tabular}{l l}
		Département:          & \departement\\
		Unité d'enseignement: & \cours\\
	\end{tabular}
\end{center}


\vfill

%\begin{center}
%	\includegraphics{Images/it.png}
%\end{center}


\begin{tabular}{l l}
	\Large	Auteurs:		& \Large \nomAuteur\\
	Professeur:				& \prof\\
	%Assistant:				& \assistant\\
	%Classe:				& \classe\\
	Date:					& \today\\
\end{tabular}



%\begin{center}
%\includegraphics[height=1.5cm]{./images/logo_reds.pdf}
%\end{center}
	% enlève le numero sur la premiere page
	\thispagestyle{empty}
	\setcounter{page}{1}
	\newpage

	\input{chapitres/Introduction.tex}

\end{document}